\section{Discussion and limitations}\label{sec:disc}

\paragraph{What was learned.}
Certifying the extremes of partisan districting at realistic resolution is possible, but only where the regime supplies structure that a relaxation can exploit. Connectivity is a local property, so a coarse quotient graph cannot see it (Proposition~\ref{prop:blind} and Section~\ref{sec:negative}); the subdivision budget is what turns the county level into an informative abstraction. Counting the budget by \emph{pieces} matters: allowing one county to be cut among all districts for a single ``split'' leaves a large unconstrained sub-problem inside that county and makes the relaxations of the states with $k\ge3$ far looser. With $k=2$ every regime of this type coincides and the loop closes essentially all brackets to half a point of district share; with $k\ge3$ most brackets remain several points wide near the margin at which the relaxation stops being spuriously feasible, and each such query then requires reasoning about the exact boundary of a single district at precinct level, which is the problem the abstraction was meant to avoid. We report those widths rather than tuning the time limits per state.

\paragraph{Reading the numbers.}
$\sigma^*(1)$ minus a party's statewide share measures the ``packing gain'' the boundaries can deliver to that party's single best district; $\sigma^*(k)$ is pinned by the statewide share (the turnout-weighted mean of the district margins equals the statewide margin), and the seats in between interpolate. Because the frontier is a capacity, its exact value at, say, 55\% is a property of the vote geography and the regime, not a forecast: it says nothing about whether any map with that many safe districts would or could be enacted. Both parties are always computed with identical code and, for both, many entries are ``zero'': a party with 39\% statewide (Rhode Island Republicans) provably cannot make a majority district in any plan of the regime.

\paragraph{Limitations.}
\begin{itemize}[leftmargin=*]
\item \emph{Atoms.} The plan class respects precinct boundaries. Splitting precincts enlarges the plan class, so our upper bounds do not bound such plans; the resolution ladder of Section~\ref{sec:robust} indicates how quickly the frontier saturates as atoms are refined, but with observed (not allocated) votes a block-level statement would require an allocation model.
\item \emph{Regime.} A tolerance of $\pm1\%$ and a budget of $k-1$ county splits are modelling choices, not legal standards; the paper makes no legal claim. Compactness, minority-opportunity and other legal constraints are not modelled. They only shrink the plan class, so by Proposition~\ref{prop:mono} our certified upper bounds remain valid upper bounds for any such sub-regime (the lower-bound plans need not satisfy the extra rules). Counties are the subdivisions everywhere, including New England.
\item \emph{Data.} Shares are two-party shares from the 2020 presidential race (robustness variants use every available statewide contest); third parties, turnout changes and other elections are outside the model, and a handful of precincts straddle counties. Island states (Hawaii) are excluded because contiguity across water is ambiguous. Artificial bridge edges (Rhode Island: one; Arkansas: four) are documented in Table~\ref{tab:instances}.
\item \emph{Solvers.} Upper bounds rest on an exact-integer CP-SAT search. We audit them against an independent MILP encoding and brute-force enumeration, but no machine-checkable proof format is produced; certificates are reproducible rather than independently verifiable objects. Lower-bound plans are stored and independently verified.
\item \emph{Coverage.} Only states with $2\le k\le5$ were run (16 states; Hawaii is excluded as above), and time limits (Section~\ref{sec:results}) leave some brackets open. The six single-district states are trivial ($F(m)=1$ exactly when the statewide share is at least $\tfrac12+m$); the 27 states with $k\ge6$ were not attempted, and nothing here says how the method scales to them.
\end{itemize}

\paragraph{Extensions.}
The hierarchy is agnostic to the atom: block-level atoms with an explicit downscaling model, or an outer relaxation over all downscalings, would extend the certificates below the precinct. Additional necessary conditions (compactness measures, opportunity districts) can be added to the relaxation, and the separability of Proposition~\ref{prop:sep} turns state certificates into national ones once the larger states are within reach. The negative result suggests where the algorithmic effort should go: strengthening the relaxation \emph{inside} a split county (for example with connectivity-aware local hulls, which we tried in simple form without gain on New Hampshire) and stronger primal search near the transition margin.

\paragraph{Use.}
The frontier quantifies the mathematical capacity of district boundaries under a stated regime. It is symmetric in the parties, recommends no map and gives no evidence about intent. We hope certified brackets make claims about what is or is not achievable by redistricting easier to check.
